\subsection{Telescopic Cascode Op-Amp}\label{sec:amp_telescopic}

\subsubsection{Overview}
A differential pair \(M_1\), \(M_2\) cascoded by \(M_3\), \(M_4\) (NMOS) and
loaded by PMOS cascode current sources \(M_5\)--\(M_8\), all in one stack
between \(V_{dd}\) and the tail \(M_9\). Each half is the cascode amplifier
of Section~\ref{sec:amp_cascode} with a cascoded load. Highest gain and
speed per unit power of single stage op-amps, at the cost of swing.

\subsubsection{Schematic (half circuit)}
\begin{circuitfig}
  \draw (0,0) node[nmos](M1){}; \node[right] at (0.3,0) {$M_1$};
  \draw (0,2) node[nmos](M3){}; \node[right] at (0.3,2) {$M_3$};
  \draw (0,4.3) node[pmos](M5){}; \node[right] at (0.3,4.3) {$M_5$};
  \draw (0,6.3) node[pmos](M7){}; \node[right] at (0.3,6.3) {$M_7$};
  \draw (M1.G) -- ++(-0.5,0) node[circ]{} node[left]{$\frac{V_{id}}{2}$};
  \draw (M3.G) -- ++(-0.5,0) node[left]{$V_{B3}$};
  \draw (M5.G) -- ++(-0.5,0) node[left]{$V_{B5}$};
  \draw (M7.G) -- ++(-0.5,0) node[left]{$V_{B7}$};
  \draw (M1.S) -- ++(0,-0.3) node[ground]{} node[right=2mm]{(virtual ground)};
  \draw (M1.D) -- (M3.S); \draw (0,1) node[circ]{} node[left]{$V_x$};
  \draw (M3.D) -- (M5.D);
  \draw (M5.S) -- (M7.D);
  \draw (M7.S) node[vcc]{$V_{dd}$};
  \draw (0,3.15) node[circ]{} -- (1.8,3.15) node[circ]{} node[right]{$V_{o1}$};
\end{circuitfig}

\subsubsection{Transfer Function}
\begin{enumerate}
  \item \textbf{Large Signal, Time Domain, DC Input Analysis (Q5)}:
    every device carries \(\frac{I_T}{2}\) (tail \(I_T\)). Stacking five
    saturated devices from ground to \(V_{dd}\), each needing
    \(|V_{DS}| \geq |V_{ov}|\):
    \begin{align*}
      V_{O,\min} &= V_{DS9,\min} + V_{DS1,\min} + V_{DS3,\min} = V_{ov9} + V_{ov1} + V_{ov3} \\
      V_{O,\max} &= V_{dd} - |V_{DS7,\min}| - |V_{DS5,\min}| = V_{dd} - |V_{ov7}| - |V_{ov5}| \\
      \text{swing per side} &= V_{O,\max} - V_{O,\min} = V_{dd} - V_{ov9} - V_{ov1} - V_{ov3} - |V_{ov5}| - |V_{ov7}|
    \end{align*}
    (the differential output \(V_{O1} - V_{O2}\) swings twice this). Unlike the
    folded cascode, the input common mode is tied to the cascode bias: \(M_1\)
    saturated needs \(V_X = V_{B3} - V_{GS3} \geq V_{IC} - V_{th1}\), so
    \(V_{IC} \leq V_{B3} - V_{GS3} + V_{th1}\), and \(V_{B3}\) in turn sets
    \(V_{O,\min}\).
  \item \textbf{Small Signal, Frequency Domain, AC Input Analysis (Q4)}:
    the half circuit (\(V_s = 0\), Section~\ref{sec:amp_diffpair}) is a cascode
    with load impedance \(Z_{up} = r_{o5} + r_{o7} + g_{m5}' r_{o5} r_{o7}\)
    (common gate \(M_5\) over \(R_S = r_{o7}\), Section~\ref{sec:amp_cg}). With
    \(G_L = \frac{1}{Z_{up}}\), \(G_m \approx g_{m1}\) and \(V_{in} = \frac{V_{id}}{2}\)
    in the cascode result \(A_V = -G_m(Z_{dn} \parallel R_L)\) of
    Section~\ref{sec:amp_cascode}:
    \begin{align*}
      Z_{dn} &= r_{o1} + r_{o3} + g_{m3}' r_{o1} r_{o3} \hspace{1cm}
      V_{o1} = -g_{m1}(Z_{dn} \parallel Z_{up}) \frac{V_{id}}{2}
    \end{align*}
    The other half is driven by \(-\frac{V_{id}}{2}\), so \(V_{o2} = -V_{o1}\) and
    \begin{align*}
      A_d &= \frac{V_{o1} - V_{o2}}{V_{id}} = \frac{2 V_{o1}}{V_{id}} = -g_{m1}(Z_{dn} \parallel Z_{up}) \approx -g_{m1}\left[(g_{m3} r_{o3} r_{o1}) \parallel (g_{m5} r_{o5} r_{o7})\right]
    \end{align*}
\end{enumerate}

\subsubsection{Analysis}
\begin{enumerate}
  \item \textbf{Gain}: all devices alike (\(g_m\), \(r_o\)),
    \(Z_{dn} \approx Z_{up} \approx g_m r_o^2\), so
    \(Z_{dn} \parallel Z_{up} \approx \frac{g_m r_o^2}{2}\) and
    \(A_d \approx -g_m \cdot \frac{g_m r_o^2}{2} = -\frac{(g_m r_o)^2}{2}\).
  \item \textbf{Bandwidth}: \(C_L\) at each output,
    \(V_{o1}(\frac{1}{Z_{dn} \parallel Z_{up}} + sC_L) = -g_{m1}\frac{V_{id}}{2}\) gives
    \(p_1 = \frac{1}{(Z_{dn} \parallel Z_{up}) C_L}\) and
    \(GB = \frac{|A_d|\, p_1}{2\pi} = \frac{g_{m1}(Z_{dn} \parallel Z_{up})}{2\pi (Z_{dn} \parallel Z_{up}) C_L} = \frac{g_{m1}}{2\pi C_L}\);
    the non-dominant pole is at the cascode node \(V_x\), where (once \(C_L\)
    shorts the output) the source of \(M_3\) presents \(\approx \frac{1}{g_{m3}}\):
    \(p_2 \approx \frac{g_{m3}}{C_x}\), which is high because \(M_3\) is an NMOS
    common gate.
  \item \textbf{Slew, Large Signal, Time Domain, AC Input Analysis (Q6)}: with all of \(I_T\) in \(M_1\), side 1 has \(I_T\)
    pulled down and \(\frac{I_T}{2}\) supplied by the load, a net \(\frac{I_T}{2}\)
    out of \(C_L\); side 2 has \(M_2\) off and \(\frac{I_T}{2}\) into \(C_L\):
    \[
      \frac{dV_{O1}}{dt} = -\frac{I_T}{2 C_L} \hspace{1cm}
      \frac{dV_{O2}}{dt} = +\frac{I_T}{2 C_L} \hspace{1cm}
      SR = \left|\frac{d(V_{O1} - V_{O2})}{dt}\right| = \frac{I_T}{C_L}
    \]
    i.e. \(\frac{I_T}{2 C_L}\) per side and \(\frac{I_T}{C_L}\) differentially.
  \item \textbf{Common mode feedback}: both outputs are high impedance nodes, so
    the output common mode is set by a CMFB loop adjusting \(V_{B7}\) or the tail.
\end{enumerate}

\subsubsection{Question}
\(I_T = 200\,\mu\)A, every \(|V_{ov}| = 0.2\,\)V, \(K_{CLM} = 0.1\,\)V\(^{-1}\) for all
devices, body effect neglected, \(C_L = 1\,\)pF per output, \(V_{dd} = 1.8\,\)V. Find
\(A_d\), \(GB\), \(SR\) and the swing per side.
\paragraph{Solution:}
\begin{align*}
  I_D &= \frac{I_T}{2} = 100\,\mu A \hspace{1cm} g_m = \frac{2 \cdot 100\,\mu A}{0.2\,V} = 1\,mS \hspace{1cm} r_o = \frac{1}{0.1 \cdot 100\,\mu A} = 100\,k\Omega \\
  Z_{dn} &= r_{o1} + r_{o3} + g_{m3} r_{o1} r_{o3} = 100\,k + 100\,k + 10^{-3} \cdot 10^{10} = 10.2\,M\Omega \hspace{1cm} Z_{up} = 10.2\,M\Omega \\
  Z_{dn} \parallel Z_{up} &= \frac{10.2\,M}{2} = 5.1\,M\Omega \\
  A_d &= -g_{m1}(Z_{dn} \parallel Z_{up}) = -10^{-3} \cdot 5.1 \times 10^{6} = -5100 \hspace{1cm} (20\log_{10} 5100 = 74.2\,dB) \\
  GB &= \frac{g_{m1}}{2\pi C_L} = \frac{10^{-3}}{2\pi \cdot 10^{-12}}\,Hz = 159\,MHz \hspace{1cm}
  SR = \frac{I_T}{C_L} = \frac{200\,\mu A}{1\,pF} = 200\,V/\mu s \\
  \text{swing} &= V_{dd} - V_{ov9} - V_{ov1} - V_{ov3} - |V_{ov5}| - |V_{ov7}| = 1.8 - 5 \cdot 0.2 = 0.8\,V
\end{align*}
